% Autor: Elias Welt
\section{Einstellungen und Benutzerpräferenzen (Settings.jsx)}

Die Einstellungskomponente \texttt{Settings.jsx} ermöglicht dem Nutzer die individuelle Anpassung der Anwendungsdarstellung. Sie demonstriert ein durchgängiges Muster zur persistenten Speicherung von UI-Präferenzen im Browser.


\subsection{Präferenz-Initialisierung beim Laden}

Beim Starten der Komponente werden alle Benutzereinstellungen aus dem \texttt{localStorage}~\cite{MDNLocalStorage} geladen und als initialer State gesetzt. Dieses Vorgehen stellt sicher, dass die Einstellungen über Seitenreloads und Browsersitzungen hinweg erhalten bleiben:

\begin{lstlisting}[language=JavaScript, caption={Laden gespeicherter Präferenzen}]
useEffect(() => {
    const currentTheme = localStorage.getItem('theme');
    if (currentTheme === 'light') setLightMode(true);

    const storedAnimations = localStorage.getItem('animationsEnabled');
    if (storedAnimations !== null)
        setAnimationsEnabled(storedAnimations === 'true');

    const storedCompact = localStorage.getItem('compactMode');
    if (storedCompact !== null)
        setCompactMode(storedCompact === 'true');
}, []);
\end{lstlisting}

Der Vergleich \texttt{storedValue !== null} ist dabei bewusst gewählt: Er unterscheidet zwischen einem nicht vorhandenen Eintrag (erster Aufruf -- Standard gilt) und einem explizit auf \texttt{false} gesetzten Eintrag (Nutzer hat ausgeschaltet).
\newpage

\subsection{Theme-Umschaltung}

Die Umschaltung zwischen Dark- und Light-Mode erfolgt durch direktes Manipulieren des \texttt{data-theme}-Attributs am \texttt{<html>}-Element. Dieses Attribut ist der Ankerpunkt für das CSS-Variablen-Theming:

\begin{lstlisting}[language=JavaScript, caption={Direkte DOM-Manipulation zur Theme-Umschaltung}]
const handleLightModeToggle = () => {
    const newLightMode = !lightMode;
    setLightMode(newLightMode);
    if (newLightMode) {
        document.documentElement.setAttribute('data-theme', 'light');
        localStorage.setItem('theme', 'light');
    } else {
        document.documentElement.removeAttribute('data-theme');
        localStorage.removeItem('theme');
    }
};
\end{lstlisting}

Die Entfernung des Attributs (\texttt{removeAttribute}) statt setzen auf einen leeren String ist bewusst gewählt: Ohne das Attribut greift das Standard-CSS (\texttt{:root}) -- also der Dark-Mode -- und es existiert kein leerer Attributwert im DOM.

\subsection{Schließen der Einstellungsseite}

Die Einstellungsseite ist als Modal-Route realisiert. Das Schließen erfolgt nicht durch einen hartcoded Redirect, sondern durch \texttt{navigate(-1)} -- den Zurück-Schritt in der Browser-Historie. Dadurch gelangt der Nutzer immer zur zuletzt besuchten Seite zurück, unabhängig davon, von wo aus die Einstellungen geöffnet wurden.
\newpage
